% =====================================================================
% Capítulo 5 · Diseño (borrador PARCIAL: solo lo decidido e implementado)
% Fuentes: código de apps/web, packages/cms y apps/web/supabase/schemas;
% docs/tfg/DECISIONES.md (ADR-001…016), BITACORA.md (B-04…B-09, B-14),
% TRAZABILIDAD.md, PROGRESO.md; packages/cms/AGENTS.md;
% .claude/skills/rls-review.
% Estado de las fuentes: 2026-09-29 (F2.4b cerrada).
% Paquetes necesarios: booktabs, tabularx, todonotes, tikz
%   (bibliotecas: positioning, arrows.meta, fit, backgrounds, calc).
% Alternativas en PlantUML: diagramas/arquitectura.puml,
%   diagramas/flujo-peticion-cms.puml
% =====================================================================

\chapter{Diseño}
\label{cap:diseno}

% PENDIENTE: borrador parcial. Faltan el modelo de datos completo, pagos
% (RF-07, F5), roles orientados a pymes (F5), configurabilidad (F5) y el
% diseño de las pantallas del CMS que aún no están hechas (F2.4c–F2.8).

Este capítulo describe el diseño de PymeKit en lo que ya está decidido e
implementado: la arquitectura general, el modelo multi-tenant y de control de
acceso, la estrategia de seguridad en profundidad y el proceso con el que se
ha verificado esa seguridad. Cada decisión remite al ADR en que se justifica.

% PENDIENTE (P-06): este capítulo presenta el diseño resultante. La plataforma
% se construyó partiendo de una base de código SaaS existente, con licencia, y
% de un CMS para el mismo tipo de base de datos (ADR-001). Acordar con el
% tutor cómo se indica aquí qué partes del diseño son heredadas y cuáles
% propias; como orientación (docs/tfg/MAPA-REFERENCIAS.md): el modelo de
% cuentas y el RBAC de equipos y del CMS son heredados; la integración del CMS
% en la web (ADR-011), el pegamento super-admin/CMS (ADR-014), el acceso a la
% consola (ADR-016) y las correcciones de seguridad (ADR-015) son propios.

% ---------------------------------------------------------------------
\section{Arquitectura general}
\label{sec:arquitectura}

\subsection{Monorepo de paquetes}

El código se organiza en un único \textit{monorepo} (un repositorio con
varias aplicaciones y paquetes que se versionan juntos) gestionado con
Turborepo~\cite{turborepo} y pnpm (ADR-002). Contiene:

\begin{itemize}
  \item \texttt{apps/web}: la aplicación web, que incluye tanto la
        aplicación SaaS como la consola de administración y el CMS;
  \item \texttt{apps/web/supabase}: los esquemas, migraciones, datos de
        prueba y pruebas pgTAP de \emph{toda} la base de datos;
  \item \texttt{apps/e2e}: las pruebas E2E con Playwright;
  \item \texttt{packages/*}: los paquetes reutilizables de la plataforma
        (\texttt{@pymekit/*}: interfaz, autenticación, cuentas, equipos,
        facturación, administración, i18n\ldots) y, en
        \texttt{packages/cms/*}, los del CMS (\texttt{@pymekit/cms-*}).
\end{itemize}

En la fecha de este borrador el monorepo tiene 47 paquetes, sin ciclos de
dependencias (RNF-04). Los paquetes declaran sus dependencias de forma
explícita y separan las exportaciones de cliente y de servidor: por ejemplo,
los paquetes del CMS con rutas y servicios son solo de servidor, y el código
de la interfaz del CMS vive en paquetes cliente aparte
(\texttt{@pymekit/cms-ui-core}, \texttt{@pymekit/cms-data-explorer-ui}\ldots).
Esta separación no es solo una cuestión de orden: evita que código con
credenciales de servidor acabe en el navegador y, como se comprobó en la
incidencia B-16, también evita ciclos entre paquetes. La \textit{build} de la
web incluye además un control que falla si el cliente de base de datos del
servidor aparece en un fragmento (\textit{chunk}) destinado al navegador.
% PENDIENTE: el nº de paquetes (47) procede de TRAZABILIDAD (RNF-04); volver a
% contarlo al cerrar F2.

\subsection{Pila tecnológica}

\begin{itemize}
  \item \textbf{Aplicación web:} TanStack Start~\cite{tanstack-start} con
        React 19. Proporciona enrutado por ficheros (TanStack
        Router~\cite{tanstack-router}), renderizado en el servidor
        (\textit{Server-Side Rendering}, SSR), cargadores de datos por ruta
        (\textit{loaders}) y \textit{server functions}: funciones que se
        ejecutan en el servidor y se invocan desde el cliente con tipos
        compartidos (RF-02). Las cachés de datos del cliente usan TanStack
        Query~\cite{tanstack-query}.
  \item \textbf{Datos, autenticación y almacenamiento:}
        Supabase~\cite{supabase-docs}, que ofrece PostgreSQL, un servicio de
        autenticación basado en JWT y almacenamiento de ficheros.
  \item \textbf{API del CMS:} Hono~\cite{hono} (\textit{framework} HTTP
        ligero basado en el estándar \texttt{Request}/\texttt{Response}) y
        Drizzle ORM~\cite{drizzle} para el acceso a la base de datos.
  \item \textbf{Pagos:} Stripe (RF-07). % PENDIENTE: diseño de pagos en F5.
  \item \textbf{Interfaz:} Tailwind CSS 4 y componentes accesibles propios
        del paquete \texttt{@pymekit/ui} (RNF-08).
\end{itemize}

Para que el conjunto se mantenga coherente se decidió usar una sola pila de
librerías en toda la web, CMS incluido: un único sistema de formularios, uno
de internacionalización y una sola biblioteca de componentes (ADR-013). Así
cada cosa se hace de una única forma (RNF-05) y el idioma español solo hay que
añadirlo en un sistema (RF-13).

\subsection{Una única base de datos}

La aplicación y el CMS comparten \textbf{una sola instancia} de PostgreSQL
(ADR-002). Los datos de la plataforma viven en el esquema \texttt{public} y
los del CMS (su RBAC, sus preferencias, las vistas guardadas, la auditoría y
los paneles) en un esquema propio llamado \texttt{cms} (ADR-012). Compartir la
base de datos es lo que permite que el CMS administre los datos de la propia
plataforma (RF-09) y que las mismas políticas de seguridad protejan ambas
partes.

\subsection{CMS integrado en la consola de administración}

La decisión de mayor calado del diseño es cómo integrar el CMS (ADR-011). Se
evaluaron tres opciones:

\begin{enumerate}[label=(\alph*)]
  \item \textbf{aplicaciones separadas} (la web, la interfaz del CMS y su
        API como tres servicios): la de menor esfuerzo, pero con tres
        despliegues, dos inicios de sesión y las \textit{cookies} de sesión
        duplicadas;
  \item \textbf{integración híbrida} (la API dentro de la web y la interfaz
        del CMS servida como aplicación estática): convivirían dos
        enrutadores y dos formas de programar;
  \item \textbf{integración total}: la API se monta dentro del servidor de
        la web y las pantallas se reescriben como rutas de la web.
\end{enumerate}

Se eligió la opción (c), a pesar de ser la de mayor esfuerzo estimado, porque
da un único servicio, un único inicio de sesión y un único estilo de código.
En concreto:

\begin{itemize}
  \item La aplicación Hono del CMS se monta en \texttt{/api/cms/*} mediante
        una ruta de servidor de TanStack Start (\texttt{routes/api/cms/\$.ts})
        que le reenvía la petición tal cual. El CMS comparte así proceso,
        dominio y \textit{cookies} con la web: un super-administrador que ya
        ha iniciado sesión en la consola lo usa sin volver a identificarse.
  \item Las pantallas del CMS son rutas de la web bajo
        \texttt{/admin/cms/*}, con \textit{loaders} y TanStack Query sobre un
        cliente RPC tipado de la API.
  \item Durante el SSR, los \textit{loaders} no hacen una petición de red a
        la propia API: llaman a la aplicación Hono dentro del mismo proceso,
        reenviando solo la cabecera \texttt{cookie} de la petición entrante
        (ADR-016).
\end{itemize}

La figura~\ref{fig:arquitectura} resume la arquitectura resultante.

\begin{figure}[htbp]
  \centering
  \begin{tikzpicture}[
      font=\footnotesize,
      box/.style={draw, rounded corners, align=center, minimum height=0.9cm,
                  inner sep=4pt},
      group/.style={draw, dashed, rounded corners, inner sep=8pt},
      arr/.style={-{Stealth[length=2mm]}}
    ]
    % Cliente
    \node[box, minimum width=10cm] (browser) at (0, 5.2)
      {Navegador (React 19): área privada, consola \texttt{/admin} y CMS \texttt{/admin/cms}};

    % Servidor web
    \node[box, minimum width=3.0cm] (ssr) at (-3.6, 3.3)
      {Rutas y \textit{loaders}\\(SSR, TanStack Router)};
    \node[box, minimum width=3.0cm] (sfn) at (0, 3.3)
      {\textit{Server functions}\\+ \textit{middleware}};
    \node[box, minimum width=3.0cm] (hono) at (3.6, 3.3)
      {Ruta \texttt{/api/cms/*}\\$\rightarrow$ aplicación Hono};

    \node[box, minimum width=3.0cm] (sbclient) at (-1.8, 1.7)
      {Cliente Supabase\\(sesión del usuario)};
    \node[box, minimum width=3.0cm] (drizzle) at (3.6, 1.7)
      {Servicios Drizzle\\(\textit{claims} del usuario)};

    \begin{scope}[on background layer]
      \node[group, fit=(ssr)(sfn)(hono)(sbclient)(drizzle),
            label={[font=\small\bfseries]left:\rotatebox{90}{apps/web (TanStack Start)}}] (web) {};
    \end{scope}

    % Supabase
    \node[box, minimum width=2.6cm] (auth) at (-3.6, -0.4) {Auth\\(JWT, MFA)};
    \node[box, minimum width=4.4cm] (pg) at (0.4, -0.4)
      {PostgreSQL + RLS\\esquemas \texttt{public} y \texttt{cms}};
    \node[box, minimum width=2.2cm] (st) at (4.2, -0.4) {Storage};
    \begin{scope}[on background layer]
      \node[group, fit=(auth)(pg)(st),
            label={[font=\small\bfseries]below:Supabase (una única instancia)}] (sb) {};
    \end{scope}

    % Externo
    \node[box] (stripe) at (-7.6, 1.7) {Stripe};

    \draw[arr] (browser) -- (ssr);
    \draw[arr] (browser) -- (sfn);
    \draw[arr] (browser) -- node[right, font=\scriptsize]{RPC} (hono);
    \draw[arr] (ssr) -- (sfn);
    \draw[arr, dashed] (ssr.east) ++(0,-0.25) -- node[below, font=\scriptsize]{SSR en proceso} ([yshift=-0.25cm]hono.west);
    \draw[arr] (sfn) -- (sbclient);
    \draw[arr] (hono) -- (drizzle);
    \draw[arr] (sbclient) -- (pg);
    \draw[arr] (sbclient) -- (auth);
    \draw[arr] (drizzle) -- (pg);
    \draw[arr] (stripe.east) -- node[above, font=\scriptsize]{\textit{webhooks}} (stripe.east -| web.west);
  \end{tikzpicture}
  \caption{Arquitectura general de PymeKit: una única aplicación web con el
    CMS integrado y una única instancia de Supabase.}
  \label{fig:arquitectura}
  % Texto alternativo: el navegador se comunica con la aplicación web
  % (TanStack Start), que contiene las rutas con SSR, las server functions y
  % la ruta /api/cms que monta la aplicación Hono del CMS. Las server
  % functions usan el cliente Supabase con la sesión del usuario; la API del
  % CMS usa servicios Drizzle con los claims del usuario. Ambos acceden a una
  % única instancia de Supabase (Auth, PostgreSQL con RLS en los esquemas
  % public y cms, y Storage). Stripe notifica a la web mediante webhooks.
\end{figure}
% PENDIENTE: revisar el ajuste fino de posiciones al compilar (no se ha
% podido compilar en el entorno donde se generó el borrador). Versión
% alternativa en diagramas/arquitectura.puml.
% Nota: los webhooks de Stripe entran por la ruta de servidor
% /api/billing/webhook de la web (apps/web/src/routes/api/billing/webhook.ts);
% el diseño de pagos se completa en F5.

% ---------------------------------------------------------------------
\section{Modelo multi-tenant y control de acceso}
\label{sec:multitenant-rbac}

\subsection{Cuentas personales y de equipo}

El \textit{tenant} de PymeKit es la \textbf{cuenta}. La tabla
\texttt{accounts} contiene dos clases de cuenta, distinguidas por el campo
\texttt{is\_personal\_account}:

\begin{itemize}
  \item la \textbf{cuenta personal}, una por usuario, cuyo identificador
        coincide con el del usuario en el servicio de autenticación;
  \item la \textbf{cuenta de equipo}, pensada para representar a una pyme
        cliente (RF-06), con un propietario principal, miembros
        (\texttt{accounts\_memberships}) e invitaciones.
\end{itemize}

Los datos de negocio se enlazan a una cuenta mediante una columna
\texttt{account\_id}, y las políticas de seguridad a nivel de fila deciden el
acceso en función de la relación del usuario con esa cuenta. Cada miembro
tiene un rol con un nivel jerárquico, y cada rol tiene un conjunto de
permisos por funcionalidad (\texttt{roles.manage}, \texttt{billing.manage},
\texttt{settings.manage}, \texttt{members.manage}, \texttt{invites.manage}).
Las políticas se apoyan en funciones auxiliares como
\texttt{has\_role\_on\_account} (pertenencia) y \texttt{has\_permission}
(permiso concreto). Las escrituras destructivas se condicionan al permiso, no
a la mera pertenencia.
\todo{Añadir el diagrama entidad-relación del esquema \texttt{public} y, tras
F5, los roles orientados a pymes.}

\subsection{Control de acceso del CMS}

El CMS tiene su propio control de acceso basado en roles
(\textit{Role-Based Access Control}, RBAC) en el esquema \texttt{cms}: cuentas
del CMS, roles con rango, grupos de permisos y permisos por tabla o por
almacenamiento. Conviven así dos modelos de autorización en la misma base de
datos: el de la plataforma (el super-administrador y los equipos) y el del
CMS. La decisión de diseño (ADR-014) fue conectarlos en lugar de sustituir
uno por otro:

\begin{itemize}
  \item \textbf{El super-administrador es la raíz del CMS.} Unos
        \textit{triggers} sobre los usuarios mantienen sincronizados a los
        super-administradores con un rol de sistema \emph{Root}, el de mayor
        rango, que concede todos los permisos. Al promocionar a un usuario
        recibe el acceso al CMS, una cuenta activa y ese rol; al retirarle la
        condición pierde el acceso, y su cuenta del CMS se desactiva pero se
        conserva para la auditoría. Un índice único garantiza que solo exista
        un rol raíz.
  \item \textbf{El RBAC del CMS da acceso limitado a otro personal}
        (soporte, gestión de contenidos). Se descartó dejar solo al
        super-administrador porque incumpliría RF-09 («con permisos propios»).
\end{itemize}

Ambos grupos comparten la consola de administración, pero con límites
distintos (ADR-016): la consola admite a quien es super-administrador
\emph{o} tiene acceso al CMS; las páginas de gestión de la plataforma exigen
super-administrador, y el personal del CMS que intente abrirlas se redirige
al CMS. La barra lateral solo muestra las secciones que la API permite, pero
esa ocultación es solo de interfaz: cada ruta y cada \textit{endpoint} vuelven
a comprobar el permiso.

% ---------------------------------------------------------------------
\section{Seguridad en profundidad}
\label{sec:seguridad}

El principio que guía el diseño de la seguridad es que \textbf{la base de
datos es la frontera de autorización} (RNF-02): toda comprobación en la
interfaz o en el servidor de aplicación se considera una primera barrera, útil
para dar una buena experiencia, pero nunca suficiente. Si una capa superior
falla, las políticas de la base de datos siguen impidiendo el acceso. Este
enfoque de defensa en profundidad es coherente con las recomendaciones de
verificación de seguridad de aplicaciones de OWASP~\cite{owasp-asvs}.
% VERIFICAR: si se cita ASVS, referenciar los requisitos concretos (control
% de acceso, V4 en la versión 4.0.3) tras comprobar la versión vigente.

\subsection{RLS como frontera}

PostgreSQL permite asociar a cada tabla políticas de seguridad a nivel de
fila (\textit{Row-Level Security}, RLS)~\cite{postgresql-rls}: predicados que
el propio motor añade a cada consulta según el usuario que la ejecuta.
Supabase~\cite{supabase-rls} expone la base de datos a clientes
autenticados y traslada la identidad del usuario (su JWT) a la sesión de
PostgreSQL, de modo que las políticas pueden decidir con ella. En PymeKit:

\begin{itemize}
  \item toda tabla tiene RLS activado; se revocan todos los privilegios por
        defecto a los roles de la API y se conceden solo los verbos que la
        funcionalidad usa (y el \texttt{UPDATE}, por columnas, para que no se
        pueda mover una fila a otra cuenta);
  \item las políticas de \texttt{UPDATE} repiten en su cláusula
        \texttt{WITH CHECK} las comprobaciones del \texttt{INSERT}, porque sin
        ella PostgreSQL solo valida la fila antigua~\cite{postgresql-createpolicy}
        (véase B-05);
  \item las funciones \texttt{security definer}, que se ejecutan con
        privilegios elevados y se saltan RLS, deben aplicar sus propias
        comprobaciones de acceso y justificarlas en un comentario;
  \item el cliente de servidor que ignora RLS se reserva para casos
        excepcionales, tras validar a mano los permisos.
\end{itemize}

\subsection{MFA obligatorio en la consola y en el CMS}

El servicio de autenticación distingue el nivel de garantía de la sesión
(\textit{Authenticator Assurance Level}, AAL)~\cite{supabase-mfa}:
\texttt{aal1} tras la contraseña y \texttt{aal2} tras completar el segundo
factor. El diseño exige \texttt{aal2} en todas las capas que dan acceso a la
administración:

\begin{itemize}
  \item en la base de datos, la función que identifica al super-administrador
        devuelve falso si la sesión no es \texttt{aal2}; unas políticas
        restrictivas exigen \texttt{aal2} a quien tiene un factor configurado;
        y el acceso al CMS lo decide la función
        \texttt{cms.verify\_admin\_access()}, que comprueba la cuenta activa
        del CMS y el MFA;
  \item en la interfaz, la consola \texttt{/admin} exige siempre
        \texttt{aal2}: con un factor configurado redirige a la verificación y
        vuelve después a la página pedida; sin factor responde 404 (ADR-016).
\end{itemize}

En PymeKit el MFA es \textbf{obligatorio por defecto} en el CMS: una opción
de configuración ausente o no válida se interpreta como «obligatorio».

\subsection{Middleware más comprobación en la base de datos}

La figura~\ref{fig:flujo-cms} muestra el recorrido de una petición a la API
del CMS y las comprobaciones que encuentra en cada capa:

\begin{enumerate}
  \item \textbf{Guardas de ruta} de la consola (\texttt{/admin} y
        \texttt{/admin/cms}): sesión \texttt{aal2} y acceso de administración.
  \item \textbf{Ruta \texttt{/api/cms/\$}} de TanStack Start: reenvía la
        petición a Hono. El código del CMS se importa dinámicamente dentro
        del manejador para que nunca llegue al navegador.
  \item \textbf{Hono}: cabeceras de seguridad, protección CSRF propia (la
        global de la web solo cubre las \textit{server functions}) y
        \textit{middleware} de autenticación, que verifica el JWT (401 sin
        sesión), exige el \textit{claim} de acceso al CMS (403) y que el
        usuario no esté bloqueado.
  \item \textbf{Servicio Drizzle}: cada transacción fija los
        \textit{claims} del usuario y el rol \texttt{authenticated} en la
        sesión de PostgreSQL mediante parámetros enlazados (nunca se
        interpola un valor del JWT en el SQL), de modo que las políticas RLS
        del esquema \texttt{cms} se aplican igual que a cualquier cliente.
  \item \textbf{PostgreSQL}: \texttt{cms.verify\_admin\_access()} y las
        políticas RLS toman la decisión final. Un super-administrador con
        sesión \texttt{aal1} supera el \textit{middleware}, pero no ve datos.
\end{enumerate}

\begin{figure}[htbp]
  \centering
  \begin{tikzpicture}[
      font=\footnotesize,
      step/.style={draw, rounded corners, align=center, minimum width=3.3cm,
                   minimum height=1.0cm, inner sep=3pt},
      check/.style={align=left, font=\scriptsize, text width=6.4cm},
      arr/.style={-{Stealth[length=2mm]}}
    ]
    \node[step] (s1) at (0, 0)     {Navegador\\(consola \texttt{/admin/cms})};
    \node[step] (s2) at (0, -1.6)  {Ruta de servidor\\\texttt{/api/cms/\$}};
    \node[step] (s3) at (0, -3.2)  {Aplicación Hono\\(\textit{middleware})};
    \node[step] (s4) at (0, -4.8)  {Servicio Drizzle\\(transacción)};
    \node[step] (s5) at (0, -6.4)  {PostgreSQL\\(esquema \texttt{cms})};

    \draw[arr] (s1) -- (s2);
    \draw[arr] (s2) -- (s3);
    \draw[arr] (s3) -- (s4);
    \draw[arr] (s4) -- (s5);

    \node[check, right=0.6cm of s1] {Guarda de ruta: sesión \texttt{aal2} y
      acceso de administración (si no: verificación MFA o 404)};
    \node[check, right=0.6cm of s2] {Reenvío a Hono; código del CMS importado
      solo en el servidor};
    \node[check, right=0.6cm of s3] {Cabeceras de seguridad, CSRF; JWT válido
      (401), \textit{claim} de acceso al CMS (403), usuario no bloqueado};
    \node[check, right=0.6cm of s4] {\textit{Claims} del usuario y rol
      \texttt{authenticated} fijados con parámetros enlazados};
    \node[check, right=0.6cm of s5] {\texttt{cms.verify\_admin\_access()}
      (cuenta activa, MFA) y políticas RLS: \textbf{decisión final}};
  \end{tikzpicture}
  \caption{Flujo de una petición a la API del CMS y comprobaciones de
    seguridad en cada capa.}
  \label{fig:flujo-cms}
  % Texto alternativo: diagrama vertical de cinco pasos (navegador, ruta de
  % servidor /api/cms/$, aplicación Hono, servicio Drizzle, PostgreSQL). A la
  % derecha de cada paso se indican sus comprobaciones; la última, las
  % políticas RLS y cms.verify_admin_access(), es la que toma la decisión.
\end{figure}
% Versión alternativa como diagrama de secuencia: diagramas/flujo-peticion-cms.puml

Las \textit{server functions} de la aplicación siguen el mismo patrón: una
cadena de \textit{middleware} tipado (autenticación, pertenencia al equipo,
super-administrador) valida la petición y después el cliente de Supabase con
la sesión del usuario deja que RLS decida.

\subsection{Fallo en cerrado}

Todas las comprobaciones de seguridad deben \textbf{fallar en cerrado}
(\textit{fail closed}): ante un dato ausente, ilegible o no válido, se
deniega. El diseño lo aplica, por ejemplo, a la opción que indica si el CMS
exige MFA, que se lee con una función de privilegios elevados precisamente
para que RLS no pueda ocultársela a quien tiene que decidir. Esta regla se
incorporó tras la incidencia B-06 (sección~\ref{sec:verificacion-seguridad}).

\subsection{Seguridad del entorno de desarrollo}

La seguridad también alcanza al entorno de trabajo. Los servicios locales de
la base de datos se aíslan de Internet con una regla de cortafuegos
específica para contenedores y se accede a la aplicación mediante un túnel
SSH (ADR-010, tras la incidencia B-04). Cualquier servicio nuevo en
contenedores se revisa con ese mismo criterio.

% ---------------------------------------------------------------------
\section{Decisiones de diseño principales}
\label{sec:decisiones}

La tabla~\ref{tab:decisiones} resume las decisiones de diseño que se han
tomado hasta la fecha, con las alternativas descartadas.

\begin{table}[htbp]
  \centering
  \footnotesize
  \caption{Resumen de decisiones de diseño (ADR) y alternativas descartadas.}
  \label{tab:decisiones}
  \begin{tabularx}{\textwidth}{@{}l>{\raggedright\arraybackslash}p{4.3cm}>{\raggedright\arraybackslash}Xl@{}}
    \toprule
    \textbf{ADR} & \textbf{Decisión} & \textbf{Alternativas descartadas (motivo)} & \textbf{Req.} \\
    \midrule
    001 & Partir de una base de código SaaS existente con licencia % PENDIENTE (P-06)
      & Desarrollo propio completo (inviable en plazo); otros \textit{boilerplates} abiertos (menos completos en multi-tenant y RLS); CMS genéricos como Strapi o Directus (otra base de datos u otro \textit{backend}) & RNF-01, 04 \\
    002 & Monorepo único y una sola base de datos & Dos repositorios con la misma BD (CI, despliegue y documentación duplicados; peor trazabilidad) & RF-09, RNF-04 \\
    004 & Conservar el CMS de datos y retirar el de contenidos editoriales & Mantener dos CMS (complejidad sin aportar al objetivo) & RF-09 \\
    010 & Aislar de Internet los servicios locales & Confiar en el cortafuegos del sistema (los contenedores se lo saltan, B-04) & RNF-02 \\
    011 & CMS integrado en la consola: API Hono en \texttt{/api/cms}, pantallas como rutas de la web & Apps separadas (tres servicios, dos inicios de sesión); híbrida (dos enrutadores) & RF-08--11 \\
    012 & Esquema SQL del CMS llamado \texttt{cms} & Conservar el nombre original (excepción permanente en el control de desmarcado) & RF-09 \\
    013 & Una sola pila de librerías en toda la web & Mantener las del CMS (dos formas de hacer lo mismo; dos sistemas de i18n) & RNF-05 \\
    014 & Super-admin como raíz del CMS; RBAC del CMS para el resto del personal & Solo super-admin, sin RBAC granular (incumple RF-09) & RF-08, 09 \\
    015 & Endurecer la seguridad del CMS de partida & Aceptar el código heredado sin auditar (brechas confirmadas, B-05--B-09) & RNF-02, 03 \\
    016 & Consola para super-admin y personal del CMS, siempre con \texttt{aal2} & Dejar entrar con el \textit{claim} de acceso sin comprobar el AAL (B-14) & RF-08, RNF-02 \\
    \bottomrule
  \end{tabularx}
\end{table}

% ---------------------------------------------------------------------
\section{Verificación de la seguridad}
\label{sec:verificacion-seguridad}

\subsection{Proceso de revisión adversarial}

Todo cambio que afecta a la base de datos (una política, un permiso, una
función \texttt{security definer}, una vista o una migración) se somete a una
revisión adversarial específica de RLS, definida como un procedimiento del
\textit{harness} (sección~\ref{sec:metodologia-ia}). Su premisa es que leer
las políticas no basta: hay que \textbf{demostrar} el aislamiento. Consta de
estas etapas:

\begin{enumerate}
  \setcounter{enumi}{-1}
  \item \textbf{Alcance}: qué tablas, funciones y esquemas se revisan.
  \item \textbf{Inventario de aislamiento}: políticas, permisos concedidos,
        funciones de privilegios elevados, vistas y restricciones de MFA y
        super-administrador de cada objeto.
  \item \textbf{Auditoría estática} con un catálogo de ataques (exposición de
        tablas, huecos en los predicados, capas restrictivas, funciones que
        se saltan RLS, vistas, almacenamiento y el esquema del CMS). Cada
        acierto se formula como un ataque concreto: qué usuario, qué
        sentencia y qué fila.
  \item \textbf{Demostración con pgTAP}~\cite{pgtap}: para cada hallazgo
        plausible, y para cada garantía que se quiere confirmar, se escribe
        una prueba que ejecuta el ataque \emph{como el usuario atacante} y
        comprueba que falla. Las pruebas quedan en el repositorio como
        pruebas de regresión.
  \item \textbf{Refutación independiente}: el contexto que produjo un
        hallazgo no puede confirmarlo. Un agente distinto recibe solo la
        política, el resultado de la prueba y la afirmación, y su única tarea
        es refutarla (encontrar la capa que en realidad protege o, ante un
        veredicto de «seguro», el tipo de sentencia, rol o nivel de sesión
        que no se ha probado). Una brecha solo se acepta si el refutador
        también la reproduce; un aprobado, solo si no encuentra un vector
        abierto.
  \item \textbf{Informe} con el veredicto, los hallazgos, la matriz de
        aislamiento y las correcciones.
\end{enumerate}

\subsection{Hallazgos corregidos}

La aplicación de este proceso a la base de datos del CMS en el incremento
F2.1 dio los resultados de la tabla~\ref{tab:hallazgos}, que también recoge
dos incidencias de seguridad detectadas por otras vías. Los fallos de origen
\emph{heredado} ya estaban en el código de partida; los \emph{propios} se
introdujeron durante el proyecto.

\begin{table}[htbp]
  \centering
  \footnotesize
  \caption{Incidencias de seguridad detectadas y corregidas (bitácora del proyecto).}
  \label{tab:hallazgos}
  \begin{tabularx}{\textwidth}{@{}l>{\raggedright\arraybackslash}Xll>{\raggedright\arraybackslash}p{3.1cm}@{}}
    \toprule
    \textbf{ID} & \textbf{Fallo} & \textbf{Origen} & \textbf{Grav.} & \textbf{Detección} \\
    \midrule
    B-04 & Base de datos local accesible desde Internet (los contenedores se saltaban el cortafuegos) & entorno & Alta & Investigación de un problema de acceso \\
    B-05 & Escalada de privilegios en el RBAC del CMS: \texttt{UPDATE} sin \texttt{WITH CHECK} & heredado & Alta & Auditoría, confirmada con prueba ejecutada \\
    B-06 & La comprobación de MFA del CMS fallaba en abierto & heredado & Alta & Una prueba nueva falló de forma inesperada \\
    B-07 & Lectura de la tabla de usuarios de autenticación mediante una función de consulta & heredado & Alta & Auditoría \\
    B-08 & Lectura del catálogo de roles de PostgreSQL (\texttt{pg\_authid}) & heredado & Alta & Refutación independiente \\
    B-09 & Política RLS con una condición tautológica & heredado & Baja & Auditoría \\
    B-14 & La consola \texttt{/admin} se cargaba sin segundo factor & propio & Media & Prueba manual del autor \\
    \bottomrule
  \end{tabularx}
\end{table}
% PENDIENTE: la bitácora no indica expresamente cómo se detectaron B-07 y
% B-09 (se asume «auditoría», dentro de la /rls-review de F2.1). Confirmar.

A continuación se resume cada uno.

\begin{description}
  \item[B-05 · Escalada por \texttt{UPDATE} sin \texttt{WITH CHECK}.] Las
        políticas de actualización de tres tablas de asignación de permisos
        solo tenían la cláusula \texttt{USING}. Un administrador delegado
        podía cambiar el permiso de una fila existente y colgar de un rol
        inferior un permiso que él mismo no tenía. Se añadió \texttt{WITH
        CHECK} con la misma comprobación que el \texttt{INSERT}. Lección: en
        PostgreSQL, un \texttt{UPDATE} sin \texttt{WITH CHECK} valida la fila
        antigua, no la nueva.
  \item[B-06 · MFA que fallaba en abierto.] La función de control leía la
        opción «MFA obligatorio» con los permisos del usuario, y una política
        restrictiva se la ocultaba precisamente a quien tenía el factor
        configurado pero había entrado sin él. La función veía la opción
        vacía, la interpretaba como «opcional» y dejaba pasar. Se corrigió
        leyendo la opción con una función de privilegios elevados y tratando
        su ausencia como «obligatorio». Al corregirlo falló una prueba
        heredada cuyo nombre decía «MFA obligatorio por defecto» pero cuya
        aserción comprobaba que se concedía el acceso (B-11).
  \item[B-07 · Lectura de usuarios de autenticación.] Dos funciones de
        consulta que se saltan RLS no aplicaban la validación de esquemas
        protegidos ni, una de ellas, la de acceso de administración, que sí
        aplicaban sus funciones hermanas de escritura. Con un permiso
        comodín podían devolver los usuarios con sus \textit{hashes} de
        contraseña. Ahora exigen ambas comprobaciones.
  \item[B-08 · Lectura del catálogo del sistema.] Tras corregir B-07, la
        etapa de refutación encontró que la lista de esquemas protegidos no
        incluía el catálogo de PostgreSQL, lo que permitía leer los
        \textit{hashes} de los roles de la base de datos. Se bloquea ahora
        cualquier esquema de sistema. Es la evidencia más clara del valor de
        la refutación independiente: encontró algo que la primera auditoría
        había dado por bueno.
  \item[B-09 · Condición tautológica.] En una subconsulta dentro de una
        política, una columna sin cualificar se resolvía contra la tabla de
        la subconsulta y se comparaba consigo misma, de modo que cualquier
        usuario con un rol veía todas las asignaciones. Se cualificó la
        columna.
  \item[B-14 · Consola sin segundo factor.] Al abrir la consola al personal
        del CMS, la guarda comprobaba el \textit{claim} de acceso, que
        también tiene un super-administrador que solo ha introducido la
        contraseña. La API y la base de datos seguían exigiendo MFA, así que
        no se exponían datos, pero la consola se cargaba. Lo detectó el autor
        con una prueba exploratoria; se corrigió exigiendo \texttt{aal2} en
        la guarda y se añadió una prueba E2E de regresión. Es también un
        ejemplo de la defensa en profundidad funcionando: el fallo de una
        capa no expuso datos porque las inferiores lo impidieron.
\end{description}

Las correcciones de la base de datos se recogen en ADR-015 y en tres
migraciones, cada una con sus pruebas de regresión en pgTAP. ADR-015 enumera
ocho puntos: dos brechas confirmadas (B-05 y B-08), cinco endurecimientos
(el fallo en abierto de la comprobación de MFA, B-06; la lectura de esquemas
protegidos, B-07; el acceso a los paneles; la unicidad de los marcadores del
rol raíz, y la política tautológica, B-09) y la alineación de cuatro funciones
cuyo esquema declarativo no coincidía con las migraciones (B-10).
Tras la corrección, la batería de pruebas de base de datos quedó en 57
ficheros y 1.550 pruebas pgTAP en verde.
% Fuente: PROGRESO.md (F2.1). Actualizar al cierre de F2/F6.

El mismo ADR deja registradas, para revisarlas en F2.7, cuatro debilidades
menores que hoy no constituyen una brecha activa, como que un permiso de
almacenamiento sin contenedor indicado equivalga a un comodín o que el
personal pueda insertar entradas de auditoría a su nombre.
\todo{Actualizar este párrafo con el resultado de la revisión en F2.7.}

\todo{En el capítulo de Pruebas, presentar la matriz de aislamiento por tabla
del informe de la revisión y las cifras finales de pruebas (pgTAP, unitarias
y E2E) tras F6.}
